% ED26-BM-RAT2021-10-2: Schmid, MIT 2021, complete chapter 10.2.
\chapter{Collective Rationality: Cooperation, Coordination and Group Agency}
\label{sec:p17-rationality-10-2}

Asking whether a group is rational can ask two different questions. Do the
participants' choices fit together to achieve a mutually good outcome? Or
does the group have beliefs, policies and intentions that fit together as
the commitments of one agent? Schmid's chapter distinguishes these questions
and proposes a connection: a capacity for thinking and acting from a shared
point of view \parencite[pp.~591--599]{Schmid2021CollectiveRationality}.
Neither question is answered merely by specifying a vote-counting rule.

The first use of \emph{collective rationality} concerns relations between
individual choices. Its characteristic standard is effectiveness or
efficiency, introduced through cooperation in a prisoner's dilemma.
The second concerns a group agent's attitudes across an agenda. Its
characteristic standard is consistency, in analogy with individual
rationality. The two can come apart: an effective plan can involve
inconsistent beliefs, while a consistent institution can pursue a bad goal
or make an ineffective choice. This chapter keeps these standards separate
before reconstructing Schmid's proposed bridge.
\index{collective rationality}\index{group agency}
\index{team reasoning}\index{corporate rationality}

\section{Evaluands, standards and evidence}
\label{sec:p17-rat102-standards}

\begin{hbdefinition}{def:p17-rat102-two-standards}{Participant cooperation and corporate consistency}
In assessing \emph{participant cooperation or coordination}, the evaluand
is a joint profile of individual actions. One can compare its outcomes for
the participants without assuming an additional collective agent. In the
source's cooperation example, collective success consists in securing mutual
cooperation rather than mutual defection.

In assessing \emph{corporate consistency}, the evaluand is the standpoint
attributed to a group: its beliefs, preferences, policies or intentions.
The group is assumed to have commitments across related questions. Its
agenda, acceptance procedure and the logical or practical constraints
connecting those commitments must be specified. Consistency of each
member's private attitudes is not consistency of the collective standpoint.
\end{hbdefinition}

This distinction does not require every collection to possess a mind.
For the first evaluand, interdependent choices are enough. For the second,
treating a party, university or business as an epistemic or practical agent
is a substantive premise, disputed in the source
\parencite[pp.~591--592, 594--596]{Schmid2021CollectiveRationality}.

A \emph{formal claim} follows from a specified game, preference profile or
logical constraint. A \emph{normative claim} concerns reasons, a requirement's
authority or a decision's legitimacy. A \emph{descriptive claim} concerns
actual behavior. Equilibrium calculation is not observation that people
play that equilibrium; neither by itself establishes a duty to comply with
a group. The distinction between a standard and its authority is developed
in Definition~\ref{def:p11-reasons-normative-senses}; structural relations
among attitudes are treated in Definition~\ref{def:p11-reasons-structural}.

Schmid's remarks about conventions, smaller groups' deliberation and ordinary
individual inconsistency are qualitative claims in his argument. The chapter
reports no participant cohort, experimental task, design, measured effect
or uncertainty estimate. Those remarks are not independently established
population frequencies here, and no normative conclusion is deduced from
them. His citation of Butler's 2012 discussion does not supply a qualified
original-study estimate in this treatment.

\section{Cooperation: individual improvement and joint efficiency}
\label{sec:p17-rat102-cooperation}

The prisoner's dilemma contrasts unilateral strategic improvement with a
mutually better outcome. The payoff ranking need not express selfish
monetary interests: utilities can represent other values. Incorporating
altruistic concerns does not by itself establish that the resulting ranking
ceases to have the dilemma's structure
\parencite[p.~592]{Schmid2021CollectiveRationality}.

\begin{hbdefinition}{def:p17-rat102-pd}{An exact symmetric prisoner's-dilemma model}
Two players choose \(c\) or \(d\), with supplied cardinal utilities
\[
 \begin{array}{c|cc}
      &c&d\\ \hline
 c&(R,R)&(S,T)\\
 d&(T,S)&(P,P)
 \end{array}
 \qquad T>R>P>S.
\]
The first coordinate is the row player's utility. The letters denote
payoff levels only in this example. Use independent mixtures and expected
utility as in Definition~\ref{def:p17-finite-strategic-game}.
A pure outcome is \emph{Pareto efficient within the four pure outcomes}
if no other pure outcome gives both players at least as much utility and
one player strictly more. This is not the unanimity-preservation axiom
for social rankings.
\index{prisoner's dilemma}\index{Pareto efficiency!outcome comparison}
\end{hbdefinition}

\begin{hbtheorem}{thm:p17-rat102-pd}{proved=annex:p17.rat102-pd}{Dominance, equilibrium and efficiency}
In Definition~\ref{def:p17-rat102-pd}, \(d\) strictly dominates \(c\)
against every opponent mixture. The unique independent mixed Nash
equilibrium is \((d,d)\). But \((c,c)\) strictly improves both utilities
over \((d,d)\), which is the only Pareto-inefficient pure outcome.
Individual best-response reasoning and joint improvement therefore assess
different comparisons.
\end{hbtheorem}

For \((T,R,P,S)=(4,3,1,0)\), defection improves the chooser's payoff from
3 to 4 against a cooperator and from 0 to 1 against a defector. Mutual
defection returns \((1,1)\), although \((3,3)\) is jointly available.
No sum of interpersonal utility units is needed for this contrast.
Joint availability does not mean either player can choose \((c,c)\)
unilaterally.

The source locates the problem in social theory, with Hobbesian antecedents
and a modern game-theoretic formulation in the 1950s. Its public-good analogy
concerns taking a shared benefit without contributing. Nonexcludability
alone does not supply this game's inequalities, and not every public-good
problem is a two-player dilemma. Schmid argues that institutional alignment
of incentives can have limits, leaving cooperative dispositions in need of
a rational account \parencite[p.~592]{Schmid2021CollectiveRationality}.

Schmid also mentions a possible reconstruction of Plato's discussion at
\emph{Republic} 359a along related lines, explicitly tentatively. That
historical suggestion is not a claim that Plato stated this normal-form
game or its equilibrium proof.

\subsection{Rival accounts revise the standard of assessment}

Schmid reports several responses, not one established theorem. Weirich's
2009 account assesses cooperative combinations by self-support rather than
ordinary utility maximization. McMahon's 2001 account compares the outcome
of a cooperative scheme, including one's contribution and the other
contributions actually made, with the noncooperative outcome. Gauthier's
constrained-maximization proposal is related. These are rival accounts of
rational cooperation as surveyed by Schmid; none proves that the unchanged
strict inequalities above recommend cooperation as a unilateral best
response \parencite[p.~593]{Schmid2021CollectiveRationality}.

The source also reports interpretations of the dilemma as a paradox or
tragedy: people classed as strategically rational can do worse together
than cooperators. The exact result is the formal contrast already stated,
not a contradiction in Nash equilibrium. Nash tests deviations with others
fixed, as Definition~\ref{def:p17-nash-equilibrium} makes explicit.
Efficiency compares joint outcomes instead.

Schmid's conceptual objection is that if each person merely benefits from
mutual cooperation, the special \emph{collective} character of the reasoning
is still unexplained. A useful relation between two people's choices need
not constitute thinking as a team. This motivates the coordination case.

\section{Coordination: equilibrium is not equilibrium selection}
\label{sec:p17-rat102-coordination}

In a pure coordination problem, an individually appropriate action depends
on what the other chooses; there need be no dominant action. Convention or
salience can help align expectations, but the game and mutual rationality
alone need not select that convention. Schmid's Hi--Lo example sharpens
the issue: both players can prefer one coordinated outcome while individual
best-response rationality also permits the inferior one
\parencite[pp.~593--594]{Schmid2021CollectiveRationality}.

If several coordinated outcomes have no unanimously better member, a
convention can still fix expectations without being a uniquely selected
consequence of the payoff description. Schmid calls the underlying
interdependence of expectations multiple contingency. He reports that
Parsons and Luhmann used it to motivate a systems-theoretic approach to
social explanation instead of rational-agency explanation, while noting
its subsequently diminished appeal. This is the source's historical
assessment, not a measured trend or a formal consequence requiring a
systems-theoretic solution. The recent team-reasoning response stays
within an agency-oriented account
\parencite[pp.~593--594]{Schmid2021CollectiveRationality}.

\begin{hbdefinition}{def:p17-rat102-hilo}{Hi--Lo with specified cardinal mixing units}
Let \(a>b>0\), and let the two actions \(H,L\) give supplied utilities
\[
 \begin{array}{c|cc}
      &H&L\\ \hline
 H&(a,a)&(0,0)\\
 L&(0,0)&(b,b).
 \end{array}
\]
Players draw independently; \(p,q\) denote respectively their probabilities
of \(H\). Ordinal preference for \(HH\) over \(LL\) alone does not
determine mixed comparisons or mixing probabilities.
\index{coordination!Hi--Lo game}\index{equilibrium selection}
\end{hbdefinition}

\begin{hbtheorem}{thm:p17-rat102-hilo}{proved=annex:p17.rat102-hilo}{All equilibria of the specified Hi--Lo game}
The independent mixed Nash equilibria are exactly
\[
 (p,q)=(1,1),\quad(0,0),\quad
 \left(\frac{b}{a+b},\frac{b}{a+b}\right).
\]
At the interior equilibrium each player chooses \(L\) with greater
probability \(a/(a+b)\), and receives expected utility
\(ab/(a+b)<b<a\). The pure equilibrium \(HH\) strictly Pareto improves
\(LL\). Neither action strictly dominates the other.
\end{hbtheorem}

For \(a=2,b=1\), the interior equilibrium selects \(H\) with probability
\(1/3\) and gives each player expected utility \(2/3\). This formalizes
the source's observation about greater equilibrium weight on the action
associated with the worse coordinated outcome. It is not an experimental
prediction. Arbitrary increasing rescaling of ordinal values need not
preserve the mixture calculation; see
Definition~\ref{def:p17-ordinal-interpretation-boundary}.

\subsection{Team reasoning and the identity of the chooser}

In Schmid's reconstruction of Sugden and Bacharach, a team reasoner asks
which joint action is best for \emph{us}, then identifies their part in
it. A team standpoint ranking \(HH\) above all other profiles selects
\(HH\), with each player's assigned part \(H\). This differs from asking
which of my actions is best against a fixed expectation about yours.
Implementation still requires the other contribution; choosing \(H\)
alone does not guarantee \(HH\).

The team ranking is not automatically an interpersonal sum. The Hi--Lo
players' agreement supports the particular Pareto comparison. Conflicting
interests in more general cases require an account of the team's objective.
The word \emph{we} does not supply it. Schmid reports Gold and Sugden's
proposal that aversion to exploitation can change a dilemma's evaluations
into a coordination structure favoring mutual cooperation. That is a change
in preferences or framing, not a cooperative equilibrium of the unchanged
dilemma \parencite[p.~594]{Schmid2021CollectiveRationality}.

Identification cannot simply be justified as another individual strategic
choice. If it is worthwhile only when others also identify, the account
recreates interdependent expectations. Schmid invokes Anderson's
priority-of-identity argument: the chooser's identity must be settled
before asking what that chooser should select. But a wholly subpersonal
framing account is also incomplete: people sometimes reason about the
appropriate standpoint, especially in role conflicts. There is no argument
that the purely individual standpoint must always be the default. The open
issue concerns agent identity and reasons, not a repeated equilibrium
calculation \parencite[p.~594]{Schmid2021CollectiveRationality}.

\section{Corporate rationality and aggregation}
\label{sec:p17-rat102-corporation}

Corporate rationality assumes a group with an agenda and attitudes. Its
consistency requirement resembles individual structural rationality, rather
than the dilemma's conflict between strategic and efficient choice.
A party's adopted policy is not every member's private opinion. Schmid
uses List and Pettit's group-agency account to frame this more demanding
notion of collectivity
\parencite[pp.~594--596]{Schmid2021CollectiveRationality}.

He reconstructs Buchanan's individualistic objection: outputs of agreed
procedures need not express an independent collective will or goal, and
their rational unity need not be desirable. The worry is that recognizing
group agency reduces individuals to organs of a larger entity. Schmid's
historical contrast is that Weber and Schumpeter treated individualism as a
methodological precept for particular explanatory projects, not a universal
ban on collective-agent concepts. This is the source's historical
interpretation, not something inferred from a voting table
\parencite[pp.~594--595]{Schmid2021CollectiveRationality}.

The opposed argument, traced to Arrow, values coherent collective choice
as protection against paralysis: inability to settle a course rather than
an affirmative preference for inaction. Desiring this capacity does not
show that an arbitrary aggregation mechanism supplies it.

\subsection{A cycle is not Arrow's proof}

Schmid's three members rank \(x>y>z\), \(y>z>x\) and \(z>x>y\).
Pairwise majorities give \(x>y\), \(y>z\) and \(z>x\), each by two votes
to one. This is Theorem~\ref{thm:p17-social-majority}. Private
transitivity coexists with collective intransitivity; inability to make
a coherent ranking is not indifference.

His money-pump illustration has the group paying to exchange \(x\) for
\(z\), then \(z\) for \(y\), then \(y\) for \(x\), repeatedly.
An ordinal cycle does not itself specify willingness to pay, liquidity or
stable sequential comparisons. The precise conditional claim is:

\begin{hbtheorem}{thm:p17-rat102-pump}{proved=annex:p17.rat102-pump}{Cyclic paid exchanges under explicit acceptance premises}
Suppose an agent starts with alternative \(x\) and cash \(w_0\).
Fix \(\varepsilon>0\). Whenever its cash \(w\geq\varepsilon\), assume
it accepts the applicable exchange
\[
 (x,w)\mapsto(z,w-\varepsilon),\qquad
 (z,w)\mapsto(y,w-\varepsilon),\qquad
 (y,w)\mapsto(x,w-\varepsilon).
\]
For every nonnegative integer \(k\), if \(w_0\geq3k\varepsilon\),
\(k\) rounds return \(x\) and leave
\(w_0-3k\varepsilon\) cash, losing \(3k\varepsilon\).
The conclusion depends on paid acceptance and available cash, not only
on \(z>x>y>z\).
\index{money pump!sequential acceptance premises}
\end{hbtheorem}

Arrow's theorem is stronger and differently quantified than this single
counterexample. The precise version in Theorem~\ref{thm:p17-social-arrow}
assumes a finite nonempty electorate, at least three alternatives,
every profile of strict individual orderings, and a complete transitive
weak social output. Pareto unanimity preservation and independence of
irrelevant alternatives then imply a fixed dictator. Its complete proof
is Annex~\ref{annex:p17.social-arrow}. Unrestricted domain describes
allowed input profiles, not uniqueness of a strict output ranking; the
theorem initially permits social ties. The source's phrase \emph{Pareto
efficiency} in its Arrow summary concerns unanimity preservation, not
efficiency of joint outcomes
\parencite[p.~595]{Schmid2021CollectiveRationality}.

These results do not prohibit coherent procedures without qualification.
Rank scores change independence; fixed authority changes nondictatorship;
domain restrictions change unrestricted domain. The stated single-peaked
restriction in Theorem~\ref{thm:p17-single-peaked-majority} gives a positive
majority result. Schmid's survey of Sen's weaker-transitivity response and
subsequent oligarchy results is not substituted for an exact new oligarchy
theorem here.

\subsection{The party's logically connected policy votes}

Judgment aggregation combines assent to propositions, not rankings of
alternatives. Schmid's party has three members, unanimously opposed to
deficits, whose separate majorities oppose higher taxes but favor higher
defense and other spending. Each member can nonetheless support a feasible
combination \parencite[p.~596]{Schmid2021CollectiveRationality}.
The narrative requires a financing relation: the signs \emph{increase} and
\emph{reduce} do not alone imply a deficit. This exact model supplies it;
it is not a factual claim about real budgets.

\begin{hbdefinition}{def:p17-rat102-policy}{A stipulated unit-budget constraint}
Let \(t,d,o\in\{0,1\}\) indicate a tax increase, a defense spending
increase and an increase in other spending. A zero spending indicator
denotes the corresponding reduction. From a balanced baseline, stipulate
changes in revenue and the two spending categories to be
\(2t,2d-1,2o-1\). No deficit requires
\[
                   (2d-1)+(2o-1)\leq2t.
\]
For binary variables this is equivalent to \((d\land o)\Rightarrow t\).
A tax increase finances both increases; reducing one spending category
offsets an increase in the other. All members and the group retain this
same no-deficit constraint.
\end{hbdefinition}

\begin{hbtheorem}{thm:p17-rat102-policy}{proved=annex:p17.rat102-policy}{Feasible individual policies, infeasible majority policy}
The profile
\[
 \begin{array}{c|ccc}
      &t&d&o\\ \hline
 1&1&1&1\\
 2&0&1&0\\
 3&0&0&1\\ \hline
 \text{coordinate-wise majority}&0&1&1
 \end{array}
\]
has three feasible individual rows and an infeasible majority row under
Definition~\ref{def:p17-rat102-policy}.
A priority procedure that chooses \(d,o\) by majority (taking zero in
any tied vote) and then sets
\(t=d\land o\) returns a feasible row for every binary profile.
Here it returns \((1,1,1)\), changing the direct majority verdict on taxes.
\end{hbtheorem}

This is the source's tax-and-spending contrast, not the court's two-premise
doctrinal example. Definitions~\ref{def:p17-ja-agenda} and
\ref{def:p17-ja-rational} specify agendas, assent and closure for the latter
general problem. The exact separately selected List--Pettit preprint result
in Theorem~\ref{thm:p17-ja-impossibility}, with full proof in
Annex~\ref{annex:p17.ja-impossibility}, assumes a finite electorate of
at least two members,
opposing propositions including two atoms and their conjunction, all
rational profiles, anonymity, one common systematic proposition-wise rule,
and complete, weakly consistent, deductively closed outputs. That preprint,
revised 9 October 2001, is not identified with the List--Pettit 2011 book
surveyed by Schmid.

The feasible priority procedure changes a premise: taxes depend on the
other selected propositions, not only the tax votes. The priority and
authority constructions of Theorem~\ref{thm:p17-ja-authority-priority}
provide further exact alternatives. Schmid's explicitly rough impossibility
summary thus concerns simultaneous requirements, not the impossibility of
every procedure, deliberation or coherent group agency
\parencite[pp.~596--597]{Schmid2021CollectiveRationality}.

\section{Consistency, legitimacy and trust}
\label{sec:p17-rat102-legitimacy}

Schmid preserves Peter's objection to making consistency universally
authoritative: democratic legitimacy can have reasons for its binding
force without consistency itself providing that force. Reason-responsiveness,
legitimacy and structural rationality are different properties
\parencite[p.~596]{Schmid2021CollectiveRationality}.

The corporate-agent reply is narrower. A party presenting one program or an
institution relied upon as a testifier has a functional reason to keep its
commitments coherent across connected questions. Schmid reports Pettit's
credibility objection to the party: inconsistency undermines the claim to
be committed to its public purpose. He also retains Roth's challenge that
this argument already presupposes the importance of group agency; it
cannot establish that premise for its opponent. A demand for reliability
in harder cases is a proposed standard of robust trust, not a measured
probability of correctness.

\begin{hbwarning}{def:p17-rat102-truth}{Coherence, consensus and identification are not truth}
Everyone can accept a false proposition \(p\) and compatible consequences.
Unanimity can preserve consistency while producing a false collective
belief. The policy profile can instead fail feasibility despite accurate
vote counting. Logical coherence tests whether judgments can hold together,
not whether they fit the actual world. Agreement is a relation among
attitudes, not truth.

Team identification also does not establish moral acceptability of a goal,
authority over outsiders or legitimacy of every instruction. Reasons to
repair shared commitments can coexist with reasons to revise the group's
purpose, refuse a contribution or leave it.
\end{hbwarning}

\section{The proposed bridge through plural self-awareness}
\label{sec:p17-rat102-intentionality}

The positive proposal is not another fixed voting rule. Schmid questions
whether mechanical aggregation adequately explains the underlying capacity.
A group can develop a plan and allocate tasks together rather than assemble
preexisting private opinions; jointly preparing a meal is his contrast.
Voting can still be necessary, especially for mass democratic decisions.
It can emulate a capacity under difficult conditions without explaining
the capacity itself \parencite[p.~597]{Schmid2021CollectiveRationality}.

In his reconstruction of List and Pettit's behavioral alternative, members
identify with the team without an \emph{identification gap}. A member
viewing the philosophy department's inconsistency as an external
institution's problem needs an extra private reason to care. Recognizing
it as \emph{our} inconsistency makes it a participant's concern. This is
a mode of owning a shared standpoint, not a further vote on every proposition.
\index{identification gap}\index{collective intentionality}

Schmid develops this through collective intentionality: minds can be jointly
directed at a matter, goal or value, though theories disagree about how to
analyze the state. His emphasis is plural prereflective self-awareness.
It is first-personal but not singular. An individual recognizing incompatible
intentions as their own must reopen what to do. On the proposed analogy,
members recognizing shared incompatible commitments must make up their
joint mind. Coherence thereby becomes a participants' concern, not merely
an external performance score
\parencite[pp.~597--598]{Schmid2021CollectiveRationality}.

The complete connection is an argument, not an existence theorem. Efficient
coordination needs both a good joint outcome and reasoning about one's
part. Corporate consistency needs both appropriate inputs and participants
maintaining one standpoint. A plural perspective can make the common aim
and the shared inconsistency first-personal concerns. Under that account
of joint agency, team reasoning supports contribution to the plan, while
collective awareness supports repair of conflicting attitudes. The common
capacity is rational joint action, not an invariant proposition-wise
function applied to unchanged inputs.

The source retains a disanalogy with singular practical judgment. A joint
judgment that an outcome is best for us does not settle every contribution
question: a member may need assurance that others will participate.
Nevertheless, on Schmid's proposal it supplies a \emph{pro tanto} reason;
acting against it calls for explanation, although contrary reasons may
outweigh it. It is not a categorical obligation deduced from membership
\parencite[pp.~597--598]{Schmid2021CollectiveRationality}.

Recognizing a conflict does not decide which belief to abandon, goal to
prioritize or role to follow. A capacity for revision does not ensure
successful revision on every occasion. The proposal does not prove that
saying \emph{we} creates a group agent, that its judgment is true or that
its members are infallible. Nor does it satisfy every original aggregation
axiom: deliberation can change inputs, agenda or procedure, making the
formal problem different.

\section{Consultation by task and worked contrasts}
\label{sec:p17-rat102-use}

\begin{hbmethod}{def:p17-rat102-diagnostic}{Select the evaluand before selecting a procedure}
\begin{enumerate}
\item Identify what is assessed: an individual response, joint action profile
      or continuing group's standpoint.
\item Identify the standard: fixed-opponent best response, Pareto outcome
      comparison, coherence, evidential accuracy or legitimacy.
\item State connecting premises: cardinal payoffs for mixing; input domain
      and output type for rankings; agenda and consequence for judgments;
      resource constraints for plans.
\item Calculate the contrast. The dilemma separates dominance from joint
      improvement; Hi--Lo separates equilibrium from selection; policy votes
      separate faithful counting from feasibility.
\item Name what a remedy changes: utility, team aim, profile domain,
      proposition priority or deliberative inputs. It is not the unchanged
      original problem satisfying all its old axioms.
\item For member commitments, distinguish a shared judgment from an
      outsider's recommendation and retain assurance, contrary reasons and
      role conflicts.
\end{enumerate}
\index{reasoning task!assess collective rationality}
\end{hbmethod}

A successful use diagnoses the party's incompatible package, states its
financing premise and adopts a coherent package through an explicitly
changed procedure. A failed use assumes that individually consistent
votes must yield consistency, or that majority makes the package true or
legitimate. In Hi--Lo, a team already owning the shared aim can allocate
the two \(H\) contributions. A failed use assumes one's \(H\) guarantees
coordination or that the inferior equilibrium is logically inconsistent.

This is a source-specific account of collective standards and their possible
connection through joint agency. The strategic-game, social-choice and
judgment-aggregation units retain their exact statements and proofs.
Schmid's philosophical proposal connects their concerns without merging
their input types or settling every objection to group agency.
